\section{Introduction}
Magnetic actuation can apply forces and torques to a slender robot without a distal motor. Its attraction for vascular intervention is accompanied by a difficult modeling requirement: the shape, material orientation, surrounding fluid and actuator geometry jointly determine the motion. A long-term application considered here is a three-dimensional flexible intravascular robot under space-station microgravity. Removing a dominant gravitational load does not remove blood-flow forcing, wall interaction, magnetic workspace limits or rotational inertia. The present study addresses the mathematical and computational interfaces needed for this application; its evidence consists of simulations and local analysis, rather than physiological or clinical validation.

A structure-preserving update is useful only after specifying the dynamics whose structure it preserves. A canonical Hamiltonian reduction supplies one such structure on a conservative branch. Nonconservative loads generally invalidate its preservation. A time-dependent Birkhoff form can represent a smooth loaded field locally, but an approximate transported form may represent a nearby field instead. Likewise, a learned scalar energy defines a conservative predictor, while a magnetic actuator can apply only a restricted momentum force. Confusing these objects can produce a small structural defect alongside an inaccurate physical trajectory. The distinction is particularly consequential for a three-dimensional rod with configuration-dependent mass and spatially varying coil fields.

We study these interfaces on a geometric Kirchhoff reduction, with two bending directions and twist, and use a spatial Cosserat model to examine environmental transfer. The reduction retains 12 configuration variables and all derivatives of its mass matrix. A full Biot--Savart energy provides the current-to-generalized-force map. Local one-form transport, Moser charts and endpoint generators are then evaluated on declared smooth branches. Hamiltonian learning and resonance diagnostics share the ideal-bias conservative reduction; the full-coil closed loop is identified separately. The paper consequently reports connected model interfaces with explicit scope, rather than a single experiment combining every method.

The contribution has three parts. First, we give a computable local connection from this variable-inertia rod field to a transported two-form and compatible generating update, retaining representation, chart, integration and solver errors as distinct quantities. The derivation specifies endpoint transversality, initial gauge and fixed-coordinate recursion, and recovers a nonseparable fourth-order central method on the autonomous conservative branch. These are applications and explicit compatibility statements within established geometric mechanics, not claims to originate Darboux or generating-function theory. Second, we connect magnetic potential, virtual work, feedback differentiation and constrained allocation to the same geometric reduction, and quantify the failure of a local affine field on the tested coil geometry. Third, we report the limits of learning, resonance and spatial-rod transfer: good sampled field prediction need not give positive momentum curvature; a nonzero formal resonant twist is not a KAM certificate; and feasible port execution need not improve tracking.

\section{Related work and relation to the research record}
Geometrically exact rod models distinguish centerline deformation from material-frame rotation. Their variational and computational foundations are described by \citet{Antman2005,Simo1985}; \citet{Gazzola2018} develop forward and inverse soft-filament mechanics. These models justify retaining shear/stretch in a full Cosserat calculation while imposing an explicit Kirchhoff restriction in a lower-dimensional analytical model. Magnetic methods and workspace-dependent actuation are reviewed by \citet{Abbott2020}. Ferromagnetic continuum robots illustrate the physical platform \citep{Kim2019}, while confined-lumen magnetic-robot simulation and navigation motivate the parameter and implementation lineage used here \citep{Tong2025}. Our coil array and residual-acceleration prescription remain benchmark choices rather than measured hardware or microgravity physiology.

Generating functions provide canonical difference schemes \citep{Feng1989}, and the broader theory of symplectic algorithms is developed by \citet{FengQin2010,Hairer2006}. Birkhoffian schemes extend the structural question to a noncanonical two-form \citep{Su2007}. For dissipative physical dynamics, the time compatibility of that form is essential: spatial closedness alone does not establish a representation of the supplied vector field. We use an auxiliary Moser flow to compute a local chart, retaining its time derivative. The original robot flow supplies a local existence argument but is not an available numerical input. The fourth-order formula in this paper is a central generating method for a nonseparable Hamiltonian, not a separable drift--kick composition.

Hamiltonian neural networks learn an energy whose derivatives define a field \citep{Greydanus2019}. Learning-to-KAM arguments require substantially stronger norm control and hypotheses than a finite test loss \citep{Chen2022}. Classical and finitely differentiable persistence results \citep{Poschel2001,Bounemoura2020}, and discrete results for symplectic algorithms \citep{Shang1999}, make those obligations explicit. In the present circular rod, four pairs of equal bending frequencies prevent direct reuse of a nonresonant two-mode certificate. We retain the existing finite-feature and statistical implications as conditional results and report the computed resonant terms without claiming a new invariant torus.

The evidence cutoff is the Stage~5 branch at repository commit \texttt{1c73a58}. The initial 1025 research manuscript supplies the physical and mathematical intent; the Stage~4 audit identifies corrections, and Stage~5 supplies their implementations and available validation. The present paper includes no Stage~6 or later bridge results and performs no new scientific simulation or training. Existing outputs are selected by their recorded validation status, with rejected derivative calculations and failed allocations retained as distinct evidence. The accompanying manifests map equations, numerical claims and figure series to exact source files and fields. This restriction matters: the same repository also contains historical two-mode and planar experiments whose existence does not establish a three-dimensional result.
